% \section{Arquitetura de controle}
% Nesta seção, descreve-se a arquitetura de controle adotado para a operação do manipulador em ambiente de microgravidade. O sistema é composto por três malhas principais, que são: controle da atitude da base (malha externa), planejamento do manipulador com mínima reação, e controle por torque nas juntas (malha interna). Essas malhas operam de forma coordenada.

% \subsection{Controle de atitude da base}
% A malha externa de controle é responsável por controlar a orientação da base do sistema. O erro de atitude é expresso por quaternions para evitar as singularidades associadas aos ângulos de Euler.
% Neste trabalho foi adotada a convenção de Hamilton\fn{Explicado anteriormente na seção \eqref{sec:Quaternions}}, assim, o erro de quaternion é dado por:

% \begin{equation}
%     q_e \;=\; q_{ref}^{-1} \otimes q_b \;=\;
%     \begin{bmatrix}
%         q_{e} \\[2pt] \vec {\epsilon}_e
%     \end{bmatrix},
%     \label{eq:erro_quaternion1}
% \end{equation}

% onde:
% \begin{itemize}
%     \item $q_{e}$ é a parte escalar do quaternion de erro e $\vec {\epsilon}_e$ é o vetor associado;
%     \item $q_b$ é o quaternion da atitude da base;
%     \item $q_{\mathrm{ref}}$ é o quaternion de referência e $q_{ref}^{-1}$ o seu inverso;
%     \item $\vec e_q$ define o vetor de erro de atitude usado na realimentação.
% \end{itemize}

% Sendo que:
% \begin{equation}
%     \vec e_q \;=\; 2\,\mathrm{sgn}(q_{e})\,\vec {\epsilon}_e,
%     \label{eq:erro_quaternion2}
% \end{equation}

% onde o termo $\operatorname{sgn}(q_{e})$ corresponde à \emph{função sinal}, utilizada para resolver a ambiguidade associada aos quaternions, assim, expressa-se:
% \begin{equation}
% \operatorname{sgn}(x) =
% \begin{cases}
% +1, & x > 0, \\
% 0,  & x = 0, \\
% -1, & x < 0,
% \end{cases}
% \label{eq:funcao_sinal}
% \end{equation}

% Como $q_e$ e $-q_e$ representam a mesma orientação física, existe um duplo recobrimento na representação.
% O fator $\operatorname{sgn}(q_{e})$ faz com que a parte escalar do quaternion de erro seja escolhida de forma consistente ($q_{e}\ge 0$), evitando o fenômeno conhecido como \emph{unwinding}, no qual o controlador tende a forçar uma rotação desnecessária de $360^{\circ}$ para atingir a referência.
% Portanto, o vetor de erro $\vec e_q$ em \eqref{eq:erro_quaternion2} fornece uma medida proporcional ao menor ângulo de rotação necessário para alinhar a base em relação à atitude de referência.


% Caso a referência apresente velocidade angular não nula, define-se também a velocidade de erro na mesma referência, como:
% \begin{equation}
%     \vec{\omega}_e \;=\; \vec{\omega}_b \;-\; \vec{\omega}_{\mathrm{ref}},
%     \label{eq:erro_velocidade}
% \end{equation}

% em que $\vec{\omega}_b$ é a velocidade angular da base e $\vec{\omega}_{\mathrm{ref}}$ a velocidade angular de referência.

% Com base nessas definições, a lei de controle proporcional-derivativa (PD) aplicada à base é

% \begin{equation}
%     \vec{\tau}_b \;=\; -\,K_p\,\vec e_q \;-\; K_d\,\vec\omega_e,
%     \label{eq:lei_controle_PD}
% \end{equation}

% onde $K_p$ e $K_d$ são os ganhos proporcionais e derivativos, respectivamente. 
% Neste trabalho o termo integral não é utilizado, de modo a evitar o acúmulo de erro ao longo do tempo (wind--up) 
% e por não ser usualmente empregado no ambiente espacial.


% \subsection{Planejamento do manipulador com mínima reação}
% Durante a movimentação do manipulador robótico, busca-se minimizar o efeito de reação na base. Para isso, utiliza-se a projeção da tarefa secundária no espaço nulo da Jacobiana, mantendo a tarefa primária (movimento do efetuador) inalterada. O comando cinemático é dado por:
% \begin{equation}
%     \dot{\vec{q}}_j = \mathbf{J}^\# \vec{v}_E^{\text{\, ref}} + \left( \mathbf{I} - \mathbf{J}^\# \mathbf{J} \right) \vec{z},
% \end{equation}

% onde:
% \begin{itemize}
%     \item $ \vec{v}_E^{\text{\, ref}} $ é a velocidade desejada do efetuador;
%     \item $ \mathbf{J}^\# $ é a pseudoinversa ponderada da Jacobiana;
%     \item $ \vec{z} $ é um vetor no espaço nulo, utilizado para tarefas secundárias.
% \end{itemize}

% A escolha de $ \vec{z} $ pode ser feita de acordo com o objetivo desejado.
% Caso seja para minimizar a reação na base, seleciona-se $ \vec{z} $ pertencente ao espaço nulo de uma matriz de acoplamento que relaciona $ \dot{\vec{q}}_j $ ao momento angular da base $ \vec{\omega}_b $.
% Já para evitar aproximações indesejadas aos limites físicos das juntas, pode-se adicionar um termo do tipo gradiente de barreira, definido por $ \vec{z} = -\nabla \Phi(\vec{q}_j) $, onde $ \Phi(\vec{q}_j) $ é uma função que penaliza a proximidade aos limites articulares.

% \subsection{Controle por torque nas juntas}

% A malha interna é responsável pelo controle fino de cada junta do manipulador, aplicando torque para seguir a trajetória desejada. A lei de controle adotada é do tipo proporcional-derivativa (PD), aplicada individualmente para cada junta $ j $:

% \begin{equation}
% \tau_j = -k_{p,j} e_j - k_{d,j} \dot{e}_j.
% \label{eq:controle_PD_junta}
% \end{equation}

% Com:

% \begin{equation}
% e_j = q_{j,\text{ref}} - q_j, \quad \dot{e}_j = \dot{q}_{j,\text{ref}} - \dot{q}_j.
% \label{eq:erro_PD_junta}
% \end{equation}

% Onde $ q_j $ e $ \dot{q}_j $ são a posição e velocidade atuais da junta, e os termos com subíndice “ref” indicam os valores de referência.
% Adicionalmente, impõe-se saturação nos torques para respeitar os limites físicos de cada atuador:

% \begin{equation}
% |\tau_j| \leq \tau_{j,\text{max}}
% \label{eq:saturacao_torque_junta}
% \end{equation}

% É comum aplicar um filtro no termo derivativo para reduzir o impacto de ruídos de medição. As trajetórias de referência $ q_{j,\text{ref}}(t) $ devem ser ao menos $ \mathcal{C}^2 $, como polinômios de quinta ordem, a fim de garantir suavidade nas acelerações.

% \subsection{Interação entre malhas}

% A malha externa de atitude fornece o torque de estabilização $ \vec{\tau}_b $, enquanto o manipulador executa $ \vec{v}_E^{\text{ref}} $ projetado no espaço nulo da Jacobiana, de forma a evitar reações sobre a base.
% Por fim, a malha interna aplica os torques articulares $ \tau_j $, respeitando os limites operacionais e as restrições dinâmicas.
% Todas as camadas estão sujeitas a restrições físicas (como zonas de exclusão, limites de torque e velocidade) e critérios de monitoramento contínuo durante a operação.

% %==================================================
% % ------------- Nota sobre ganhos de controle (escopo)
% %==================================================
% \subsection{Nota sobre ganhos de controle}

% A lei de controle utilizada é do tipo proporcional--derivativa (PD) nas variáveis de atitude e velocidade da base e nas coordenadas de junta.
% Os valores numéricos dos ganhos $(K_p,K_d)$, bem como limites e eventuais filtros auxiliares, serão especificados no capítulo relacionado às simulações, juntamente com os cenários de simulação e as métricas de avaliação.
